\section{问题分析}\label{sec:analysis}

\subsection{总体思路}

四个问题描述的是同一个物理过程在不同设定下的表现，因此先建立一个统一的传热传质模型，再按各问的差异替换物性、几何与判据。分析的关键有以下几点。

\parahead{一维径向近似} 药材长 $\SI{25}{cm}$、直径 $\SI{4}{cm}$，长直径比 $L/(2R_0)=6.25$（长半径比 $L/R_0=12.5$），属于细长圆柱。热风从侧面包裹药材，热量与水分主要沿半径方向进出，轴向的梯度和两端的散失相对次要。因此把问题简化为沿半径 $r$ 的一维模型，只在药材中截面上求解温度和含水率，既保留主要的输运方向，又显著降低计算量。题目所说的「各处」在本文中即指半径区间 $r\in[0,R_0]$ 上的所有位置。忽略端面所带来的偏差有多大，将在第~\ref{sec:verification}~节用一个乘积解作定量估计。

\parahead{变物性与热质耦合} 药材在干燥中含水率从 $2.55$ 降到 $0.15$ 附近，跨越一个数量级，其密度、比热、导热系数和水分扩散系数都随含水率明显变化，扩散系数还随温度按 Arrhenius 形式变化。所谓耦合，是指含水率改变热物性、温度又改变水分扩散，二者相互影响。问题一给出的是一组常物性、且扩散系数只依赖含水率的特例，此时温度方程与水分方程可以分别求解；问题二至四的物性使两者相互牵制，必须把温度场与水分场作为一个耦合系统同时推进。

\parahead{收缩区域的参考坐标变换} 问题四中药材边界随失水向内移动，计算域本身在变化。若直接在物理坐标上求解，需处理移动边界及相应的输运。为此引入随物料一起收缩的相对坐标 $x=r/R(t)$，它按当前半径衡量相对位置、使表面始终对应 $x=1$，从而把动边界问题变换到固定的单位区间上求解。在仿射收缩假设下，边界移动带来的相对对流恰好相消，方程中只剩下半径 $R(t)$ 作为已知的时间函数（详见第~\ref{sec:q4}~节）。

\parahead{基于最大含水率的烘干终止判据} 由于水分自内部向表面迁移并在表面蒸发，靠近表面的区域先干、干燥前沿由表面向中心推进，只有最湿的位置也降到要求以下，整个径向域才满足烘干要求。在本题已计算的工况下，最湿的位置是药材中心，因此「各处含水率都低于 $0.15$」等价于「最大含水率低于 $0.15$」。据此把烘干终止判据落实为：在连续重构的含水率场上，求全域最大含水率随时间下降到 $0.15$ 的时刻；算法仍取全域最大值，而不把中心恒为最大值当作前提。这样得到的时间由方程本身确定，不依赖输出时的四位小数取舍。

\subsection{各问思路}

\parahead{问题一} 在常物性、扩散系数只依赖含水率的设定下求预热阶段的场。此时温度与水分可以分别求解，重点是刻画表面附近陡峭的边界层。

\parahead{问题二} 把物性换成随含水率、温度变化的形式，温度与水分耦合，求解贯穿整个干燥过程，得到全时段的两个场。

\parahead{问题三} 作为结果判定环节，不另建控制方程，直接在问题二得到的同一条数值解上追踪全域最大含水率，求其下降到 $0.15$ 的时刻作为烘干时长。

\parahead{问题四} 在问题二、三的耦合模型基础上引入另一组物性（附录 4）与随时间收缩的半径，并在参考坐标上求解，得到收缩情形下的场与烘干时长。各问均从题给初始状态起算。

四个问题之间的承接关系如图~\ref{fig:relation} 所示：问题一使用附录 2，问题二从初始时刻起统一使用附录 3，问题三基于问题二的同一数值解，问题四同时改用附录 4 物性与给定的收缩半径。

\begin{figure}[H]\centering
\includegraphics[width=\textwidth]{fig_relation.pdf}
\caption{四个问题与模型的关系：四问共用同一套径向传热传质模型；问题一为常物性预热特例，问题二为变物性完整过程；问题四在问题二模型上作物性与几何修正，问题三在问题二的结果上按最大含水率判据确定烘干时长。}
\label{fig:relation}
\end{figure}
